% GENERATED FILE. Edit canonical lessons, then run npm run generate:books.
% canonical-source-hash: fnv1a64:9d8d17e0baf04d35
% canonical-lessons: TA-C226-arampi, TA-C226-totu, TA-C226-tummu, TA-C226-vilunku, TA-C226-mel

\chapter{To begin, To touch, To sneeze, To swallow, To chew}
\label{ch:ta-to-begin-to-touch-to-sneeze-to-swallow-to-chew}
\hlchaptermodality{226}

\begin{chapteropening}
\textbf{By the end of this chapter:} \emph{I can say to begin, to touch, to sneeze, to swallow and to chew.}

Complete the last lesson of chapter 226: I can say to begin, to touch, to sneeze, to swallow and to chew.
\end{chapteropening}

\section[\texorpdfstring{\ta{ஆரம்பி} (ārampi)}{ஆரம்பி (ārampi)}]{\ta{ஆரம்பி} (ārampi) — to begin}

\label{lesson:TA-C226-arampi}



\begin{quote}
Before the new one: say the Tamil for to take off, then the Tamil for to change.
\end{quote}

\subsection*{You'll want to know: \ta{ஆரம்பி}}
\textbf{\ta{ஆரம்பி}} — \emph{ārampi} — \textquotedblleft{}to begin\textquotedblright{}.

\begin{grammarlens}[title={What you've built}]
One more word for this chapter.
\end{grammarlens}

\subsection*{Your turn}
\begin{itemize}
\raggedright
  \item \emph{Say it:} \emph{ārampi}
  \item \emph{Say it:} \emph{ārampi}, once more
  \item \emph{Say it:} say \emph{ārampi}
  \item \emph{Recall:} say the Tamil for to take off, then the Tamil for to change, then say \emph{ārampi} again
\end{itemize}

\begin{tcolorbox}[breakable,colback=teal!4,colframe=teal!35!black,title={Before you move on}]
What does \ta{ஆரம்பி} mean? (\textquotedblleft{}To begin\textquotedblright{}.) Say it once more.
\end{tcolorbox}
\section[\texorpdfstring{\ta{தொடு} (toṭu)}{தொடு (toṭu)}]{\ta{தொடு} (toṭu) — to touch}

\label{lesson:TA-C226-totu}



\begin{quote}
Before the new one: say the Tamil for to change, then the Tamil for to begin.
\end{quote}

\subsection*{You'll want to know: \ta{தொடு}}
\textbf{\ta{தொடு}} — \emph{toṭu} — \textquotedblleft{}to touch\textquotedblright{}.

\begin{grammarlens}[title={What you've built}]
One more word for this chapter.
\end{grammarlens}

\subsection*{Your turn}
\begin{itemize}
\raggedright
  \item \emph{Say it:} \emph{toṭu}
  \item \emph{Say it:} \emph{toṭu}, once more
  \item \emph{Say it:} say \emph{toṭu}
  \item \emph{Recall:} say the Tamil for to change, then the Tamil for to begin, then say \emph{toṭu} again
\end{itemize}

\begin{tcolorbox}[breakable,colback=teal!4,colframe=teal!35!black,title={Before you move on}]
What does \ta{தொடு} mean? (\textquotedblleft{}To touch\textquotedblright{}.) Say it once more.
\end{tcolorbox}
\section[\texorpdfstring{\ta{தும்மு} (tummu)}{தும்மு (tummu)}]{\ta{தும்மு} (tummu) — to sneeze}

\label{lesson:TA-C226-tummu}



\begin{quote}
Before the new one: say the Tamil for to begin, then the Tamil for to touch.
\end{quote}

\subsection*{You'll want to know: \ta{தும்மு}}
\textbf{\ta{தும்மு}} — \emph{tummu} — \textquotedblleft{}to sneeze\textquotedblright{}.

\begin{grammarlens}[title={What you've built}]
One more word for this chapter.
\end{grammarlens}

\subsection*{Your turn}
\begin{itemize}
\raggedright
  \item \emph{Say it:} \emph{tummu}
  \item \emph{Say it:} \emph{tummu}, once more
  \item \emph{Say it:} say \emph{tummu}
  \item \emph{Recall:} say the Tamil for to begin, then the Tamil for to touch, then say \emph{tummu} again
\end{itemize}

\begin{tcolorbox}[breakable,colback=teal!4,colframe=teal!35!black,title={Before you move on}]
What does \ta{தும்மு} mean? (\textquotedblleft{}To sneeze\textquotedblright{}.) Say it once more.
\end{tcolorbox}
\section[\texorpdfstring{\ta{விழுங்கு} (viḻuṅku)}{விழுங்கு (viḻuṅku)}]{\ta{விழுங்கு} (viḻuṅku) — to swallow}

\label{lesson:TA-C226-vilunku}



\begin{quote}
Before the new one: say the Tamil for to touch, then the Tamil for to sneeze.
\end{quote}

\subsection*{You'll want to know: \ta{விழுங்கு}}
\textbf{\ta{விழுங்கு}} — \emph{viḻuṅku} — \textquotedblleft{}to swallow\textquotedblright{}.

\begin{grammarlens}[title={What you've built}]
One more word for this chapter.
\end{grammarlens}

\subsection*{Your turn}
\begin{itemize}
\raggedright
  \item \emph{Say it:} \emph{viḻuṅku}
  \item \emph{Say it:} \emph{viḻuṅku}, once more
  \item \emph{Say it:} say \emph{viḻuṅku}
  \item \emph{Recall:} say the Tamil for to touch, then the Tamil for to sneeze, then say \emph{viḻuṅku} again
\end{itemize}

\begin{tcolorbox}[breakable,colback=teal!4,colframe=teal!35!black,title={Before you move on}]
What does \ta{விழுங்கு} mean? (\textquotedblleft{}To swallow\textquotedblright{}.) Say it once more.
\end{tcolorbox}
\section[\texorpdfstring{\ta{மெல்} (mel)}{மெல் (mel)}]{\ta{மெல்} (mel) — to chew}

\label{lesson:TA-C226-mel}



\begin{quote}
Before the new one: say the Tamil for to sneeze, then the Tamil for to swallow.
\end{quote}

\subsection*{You'll want to know: \ta{மெல்}}
\textbf{\ta{மெல்}} — \emph{mel} — \textquotedblleft{}to chew\textquotedblright{}.

\begin{grammarlens}[title={What you've built}]
One more word for this chapter.
\end{grammarlens}

\subsection*{Your turn}
\begin{itemize}
\raggedright
  \item \emph{Say it:} \emph{mel}
  \item \emph{Say it:} \emph{mel}, once more
  \item \emph{Say it:} say \emph{mel}
  \item \emph{Recall:} say the Tamil for to sneeze, then the Tamil for to swallow, then say \emph{mel} again
\end{itemize}

\begin{tcolorbox}[breakable,colback=teal!4,colframe=teal!35!black,title={Before you move on}]
What does \ta{மெல்} mean? (\textquotedblleft{}To chew\textquotedblright{}.) Say it once more.
\end{tcolorbox}
